%! AP Statistics — Unit 2, Lesson 2
\documentclass[10pt]{article}
\usepackage[margin=0.6in]{geometry}
\usepackage{xcolor}
\usepackage[most]{tcolorbox}
\usepackage{enumitem}
\usepackage{multicol}
\usepackage{amsmath,amssymb}
\usepackage{array}
\usepackage{tabularx}
\tcbuselibrary{breakable}

\definecolor{sky}{HTML}{EAF3FB}
\definecolor{navy}{HTML}{1F3A5F}
\definecolor{redbox}{HTML}{FCECEC}
\definecolor{greenbox}{HTML}{EDF8F0}
\definecolor{goldbox}{HTML}{FFF7E6}
\newtcolorbox{skillbox}[2][]{
    colback=#2, colframe=navy, boxrule=0.8pt, arc=2mm,
    left=2mm,right=2mm,top=1.5mm,bottom=1.5mm,
    fonttitle=\bfseries, title={#1}, halign title=center, breakable
}
\setlist[itemize]{leftmargin=1.2em,itemsep=0.35em,topsep=0.35em}
\setlist[enumerate]{leftmargin=1.4em,itemsep=0.35em,topsep=0.35em}
\newcolumntype{C}[1]{>{\centering\arraybackslash}p{#1}}
\newcolumntype{Y}{>{\centering\arraybackslash}X}

\newcommand{\CourseName}{AP Statistics}
\newcommand{\SchoolYear}{2026--2027}
\newcommand{\MeetingLength}{55 minutes}
\newcommand{\UnitNumberName}{Unit 2: Exploring Two-Variable Data \quad}
\newcommand{\LessonNumberName}{Lesson 2.2: Representing Two Categorical Variables}

\begin{document}
\begin{center}
    {\Large\bfseries \CourseName: \SchoolYear\ (\MeetingLength) \\
    {\small\bfseries \UnitNumberName \LessonNumberName}}
\end{center}

\begin{tcolorbox}[colback=sky,colframe=navy,boxrule=0.9pt,arc=2mm,
    left=2mm,right=2mm,top=1.5mm,bottom=1.5mm]
    \textbf{Primary Objective:} Students will construct and read two-way (contingency) tables and
    segmented bar charts for two categorical variables. Students will distinguish between joint,
    marginal, and conditional frequencies, and will use these displays to explore whether an
    association exists between the two variables (Big Idea: UNC).
\end{tcolorbox}

\begin{skillbox}[Priority Ideas \& Skills]{goldbox}
\begin{minipage}[t]{0.48\linewidth}
    \textbf{AP Skill 2 --- Data Analysis}
    \begin{itemize}
        \item Construct a two-way table from raw data; identify cells as joint frequencies.
        \item Compute marginal frequencies (row totals and column totals).
        \item Construct a segmented bar chart using conditional relative frequencies to compare groups visually.
    \end{itemize}
\end{minipage}
\hfill
\begin{minipage}[t]{0.48\linewidth}
    \textbf{Key Understandings}
    \begin{itemize}
        \item Marginal frequencies describe one variable at a time (same as a one-way table).
        \item Conditional frequencies hold one variable fixed and describe the distribution of the other --- these reveal association.
        \item A segmented bar chart with conditional relative frequencies makes comparisons between groups visually clear.
    \end{itemize}
\end{minipage}
\end{skillbox}

\begin{skillbox}[{Vocabulary, Concepts, \& Theorems}]{greenbox}
\begin{minipage}[t]{0.48\linewidth}
    \begin{itemize}
        \item \textbf{Two-way table (contingency table)} --- organizes counts for two categorical variables; rows = one variable, columns = the other
        \item \textbf{Joint frequency} --- count in a single cell; describes both variables simultaneously
        \item \textbf{Marginal frequency} --- row total or column total; describes one variable ignoring the other
    \end{itemize}
\end{minipage}
\hfill
\begin{minipage}[t]{0.48\linewidth}
    \begin{itemize}
        \item \textbf{Conditional relative frequency} --- relative frequency for one variable \textit{given} a specific value of the other variable
        \item \textbf{Segmented bar chart} --- each bar represents a group; bar divided into segments proportional to conditional relative frequencies
        \item \textbf{Association (categorical)} --- if conditional distributions differ across groups, the two variables are associated
    \end{itemize}
\end{minipage}
\end{skillbox}

\begin{skillbox}[Activate Prior Knowledge \& Spiral Review]{sky}
\begin{minipage}[t]{0.48\linewidth}
    \textbf{Quick Review (3 min)}\\
    Recall relative frequency tables from Lesson 1.3. Ask: ``What if I want to cross two categorical variables at once? For example, grade level AND preferred subject?'' A one-way table can't show both --- that's the motivation for today's two-way table.
\end{minipage}
\hfill
\begin{minipage}[t]{0.48\linewidth}
    \textbf{Spiral Connections}
    \begin{itemize}
        \item Lesson 1.3: relative frequency tables (one variable) $\to$ today: two variables crossed.
        \item Lesson 2.3: we compute statistics (conditional proportions, relative risk) from these tables.
        \item Unit 8: Chi-square tests formally assess whether an association in a two-way table is statistically significant.
    \end{itemize}
\end{minipage}
\end{skillbox}

\begin{skillbox}[Hook \& Lesson]{sky}
\begin{minipage}[t]{0.48\linewidth}
    \textbf{Hook (5 min)}\\
    Present a table of data: 200 students surveyed on grade level (9th/10th) and whether they have a part-time job (Yes/No). Ask: ``Do 10th graders have jobs at a higher rate than 9th graders?'' Students immediately see a raw table isn't enough --- we need conditional percents.
\end{minipage}
\hfill
\begin{minipage}[t]{0.48\linewidth}
    \textbf{Lesson Flow (15 min)}
    \begin{enumerate}
        \item Build the two-way table together; label joint and marginal cells.
        \item Compute conditional relative frequencies for each grade level (row percents).
        \item Construct a segmented bar chart from those conditional percents.
        \item Read the chart: do the segments look different? $\to$ evidence of association or no association.
    \end{enumerate}
\end{minipage}
\end{skillbox}

\begin{skillbox}[Explicit Instruction \& Active Monitoring]{sky}
\begin{minipage}[t]{0.48\linewidth}
    \textbf{Direct Instruction (10 min)}
    \begin{itemize}
        \item Walk through computing row conditional percents vs.\ column conditional percents; when to use each (condition on the explanatory variable).
        \item Segmented bar construction: each bar totals 100\%; segments colored by response variable categories.
        \item Association check: if the segmented bars look the same for all groups $\to$ no association.
    \end{itemize}
\end{minipage}
\hfill
\begin{minipage}[t]{0.48\linewidth}
    \textbf{Active Monitoring}
    \begin{itemize}
        \item Common error: computing conditional percents by dividing by the grand total (gives joint frequencies) instead of the row or column total.
        \item Probe: ``Which variable is the explanatory? That's the one you condition on (divide by its row/column total).''
        \item Check segmented bars: do they add to 100\%? Are they on the same scale?
    \end{itemize}
\end{minipage}
\end{skillbox}

\begin{skillbox}[Group Work \& Differentiation]{redbox}
\begin{minipage}[t]{0.48\linewidth}
    \textbf{Group Activity (8--10 min)}\\
    Groups receive a two-way table (political affiliation $\times$ preferred news source, $n=300$). Groups:
    \begin{enumerate}
        \item Compute and add marginal frequencies.
        \item Compute conditional relative frequencies for each political group.
        \item Construct a segmented bar chart.
        \item Write one sentence: is there evidence of an association? Justify.
    \end{enumerate}
\end{minipage}
\hfill
\begin{minipage}[t]{0.48\linewidth}
    \textbf{Differentiation}
    \begin{itemize}
        \item \textit{Support}: Provide a partially completed table with marginal totals filled in.
        \item \textit{Extension}: Investigate Simpson's Paradox with a provided dataset where the overall association reverses when data are separated by a third variable.
        \item \textit{Technology}: Use a spreadsheet to build a segmented bar chart and compare to the hand-drawn version.
    \end{itemize}
\end{minipage}
\end{skillbox}

\begin{skillbox}[Individual Work \& Assessment]{redbox}
\begin{minipage}[t]{0.48\linewidth}
    \textbf{Exit Ticket (5 min)}\\
    Given a two-way table (gender $\times$ prefers fiction or nonfiction, $n=120$), students: (1) compute conditional relative frequencies for each gender, (2) state whether there is evidence of an association between gender and reading preference, (3) justify using specific values.
\end{minipage}
\hfill
\begin{minipage}[t]{0.48\linewidth}
    \textbf{AP-Style Check}\\
    \textit{``The table shows data on transportation method and whether students are late or on time. Compute the conditional relative frequencies and describe any association you see between transportation method and tardiness.''}\\[4pt]
    Assess: correct computation, clear comparison of conditional distributions, context included.
\end{minipage}
\end{skillbox}

\begin{skillbox}[Reinforcement \& Extension]{goldbox}
\begin{minipage}[t]{0.48\linewidth}
    \textbf{Homework / Practice}
    \begin{itemize}
        \item Construct a two-way table from a raw dataset (30 individuals, 2 categorical variables).
        \item Compute marginal and conditional relative frequencies; build a segmented bar chart; describe any association.
    \end{itemize}
\end{minipage}
\hfill
\begin{minipage}[t]{0.48\linewidth}
    \textbf{Extension / Preview}
    \begin{itemize}
        \item \textit{Extension}: Research a real two-way table from a published survey (e.g., Pew Research); describe the association and suggest a possible explanation.
        \item \textit{Preview}: Lesson 2.3 --- we will quantify association using statistics like the difference in proportions and relative risk, giving us numbers to compare rather than just visual impressions.
    \end{itemize}
\end{minipage}
\end{skillbox}

\end{document}
